\subsection{Five-Species Versus Eleven-Species Gas Study\label{ch4:speciesStudy}}
    This investigation compares the index of refraction between a $5$ and $11$
    species gas. The solver used for this case was \gls{SU2} \gls{NEMO}
    \cite{Maier2021}, which is the nonequilibrium version of \gls{SU2} that is
    coupled with \gls{Mutation++} \cite{Scoggins2014} library. The
    Park's $2T$ nonequilibrium model \cite{Park1989} model is used for this
    investigation. 

    \begin{table}[H]
        \centering
        \begin{tabular}{|c|c|c|c|c|c|}
            \hline
            Case & Species & $T_{tr}\Units{K}$ & $T_{rot}\Units{K}$ &
            $T_{vib}\Units{K}$ & $P\Units{kPa}$ \\
            \hline \hline
            $1S1$\cite{Gimelshein2022} & $5$ & $15000$ & $15000$ & $300$ & $2067$
            \\ \hline
            $2S1$ & $5$ & $8000$ & $8000$ & $300$ & $1000$     \\ \hline
            $1S2$ & $11$ & $15000$ & $15000$ & $300$ & $2067$  \\ \hline
            $2S2$ & $11$ & $8000$ & $8000$ & $300$ & $1000$    \\ \hline
        \end{tabular}
        \caption{Species studied cases}\label{tab:speciesStudyCases}
    \end{table}

    Cases in \Cref{tab:speciesStudyCases} used the mass fraction
    composition for standard Air provided in
    \Cref{tab:speciesStudyMassFraction}.
    \begin{table}[H]
        \centering
        \begin{tabular}{|c|c|c|c|c|}
            \hline
            $\ce{N}\Units{\,}$ & $\ce{O}\Units{\,}$ & $\ce{NO}\Units{\,}$ & $\ce{N_2}\Units{\,}$ &
            $\ce{O_2}\Units{\,}$ \\ \hline\hline
            $0.0$ & $0.0$ & $0.0$ & $0.79$ & $0.21$ \\ \hline
        \end{tabular}
        \caption{Mass fraction composition for cases in
        \Cref{tab:speciesStudyCases}\label{tab:speciesStudyMassFraction}}
    \end{table}

    \subsubsection{Results}
    % Temperatures %
    \Cref{fig:temperaturesSpeciesStudy} illustrates the translational and
    vibrational temperatures for the cases described in
    \Cref{tab:speciesStudyCases} and \Cref{tab:speciesStudyMassFraction}. The
    temperature plots compare a five-species gas
    (\Cref{fig:temperature1S1,fig:temperature2S1}) with an eleven-species gas
    (\Cref{fig:temperature1S2,fig:temperature2S2}). In the highest-temperature
    cases (\Cref{fig:temperature1S1,fig:temperature1S2}), thermal equilibrium
    is reached at approximately $0.5\Units{ns}$. Notably, the temperatures do
    not appear to differ significantly between the five-species and
    eleven-species gases. Similarly, in the lowest-temperature cases
    (\Cref{fig:temperature2S1,fig:temperature2S2}), thermal equilibrium is
    achieved at approximately $5\Units{ns}$, and again, there is no observable
    difference between the temperatures of the five-species and eleven-species
    gases.
    \begin{figure}[H]
        \centering
        \subfigure[Case$1S1$]
        {\label{fig:temperature1S1}\includegraphics[width=0.45\textwidth]{chapter4/fiveVsElevenSpecies/1S1_temperatures}}
        \subfigure[Case$1S2$]
        {\label{fig:temperature1S2}\includegraphics[width=0.45\textwidth]{chapter4/fiveVsElevenSpecies/1S2_temperatures}}
        \subfigure[Case$2S1$]
        {\label{fig:temperature2S1}\includegraphics[width=0.45\textwidth]{chapter4/fiveVsElevenSpecies/2S1_temperatures}}
        \subfigure[Case$2S2$]
        {\label{fig:temperature2S2}\includegraphics[width=0.45\textwidth]{chapter4/fiveVsElevenSpecies/2S2_temperatures}}
        \caption{\label{fig:temperaturesSpeciesStudy}Temperature results for
        cases in \Cref{tab:speciesStudyCases}}
    \end{figure}
    % Temperatures %

    %% Mass Fraction %%
    \Cref{fig:massFractionSpeciesStudy} illustrates the mass fractions for the cases described in
    \Cref{tab:speciesStudyCases} and \Cref{tab:speciesStudyMassFraction}. In
    the highest-temperature cases ($1S1$ and $1S2$), a higher rate of
    dissociation is observed, which is expected due to the greater energy
    available in the system. For the eleven-species case at the highest
    temperature (\Cref{fig:massFraction1S2}), a significant amount of diatomic
    nitrogen ions is generated, peaking at approximately $50\Units{ns}$.
    However, the mass fraction of other ion species do not change
    significantly. In contrast, the lowest-temperature cases ($2S1$ and $2S2$)
    exhibit a lower rate of dissociation compared to the higher-temperature
    cases. In these cases, the maximum concentration of diatomic nitrogen ions
    occurs at approximately $50\Units{ns}$. Although diatomic nitrogen ions are
    formed, their quantities remain negligible compared to the neutral species.
    \begin{figure}[H]
        \centering
        \subfigure[Case$1S1$]
        {\label{fig:massFraction1S1}\includegraphics[width=0.43\textwidth]{chapter4/fiveVsElevenSpecies/1S1_massFraction}}
        \subfigure[Case$1S2$]
        {\label{fig:massFraction1S2}\includegraphics[width=0.48\textwidth]{chapter4/fiveVsElevenSpecies/1S2_massFraction}}
        \subfigure[Case$2S1$]
        {\label{fig:massFraction2S1}\includegraphics[width=0.43\textwidth]{chapter4/fiveVsElevenSpecies/2S1_massFraction}}
        \subfigure[Case$2S2$]
        {\label{fig:massFraction2S2}\includegraphics[width=0.48\textwidth]{chapter4/fiveVsElevenSpecies/2S2_massFraction}}
        \caption{\label{fig:massFractionSpeciesStudy}Mass fraction results for cases in \Cref{tab:speciesStudyCases}}
    \end{figure}
    %% Mass Fraction %%

    %% Species Gladstone-Dale %%
    \Cref{fig:speciesGladstoneDaleStudy} shows the species \gls{GD} for the
    cases described in \Cref{tab:speciesStudyCases} and
    \Cref{tab:speciesStudyMassFraction}. The eleven-species cases ($1S2$ and
    $2S2$) demonstrate that the ionic part of the \gls{GD} constants are
    primarily influenced by the formation of charged nitric oxide. However,
    these values still appear to be low compared to those of the neutral
    species. \\

    Comparing the five-species and eleven-species cases, it is evident that the
    ion contribution to the \gls{GD} is small—approximately one order of magnitude
    less for the highest-temperature case and two orders of magnitude less for the
    lowest-temperature case. Consequently, the ion contribution does not significantly impact the
    \gls{GD}. 
    \begin{figure}[H]
        \centering
        \subfigure[Case$1S1$]
        {\label{fig:speciesGladstoneDale1S1}\includegraphics[width=0.43\textwidth]{chapter4/fiveVsElevenSpecies/1S1_speciesGladstoneDale}}
        \subfigure[Case$1S2$]
        {\label{fig:speciesGladstoneDale1S2}\includegraphics[width=0.48\textwidth]{chapter4/fiveVsElevenSpecies/1S2_speciesGladstoneDale}}
        \subfigure[Case$2S1$]
        {\label{fig:speciesGladstoneDale2S1}\includegraphics[width=0.43\textwidth]{chapter4/fiveVsElevenSpecies/2S1_speciesGladstoneDale}}
        \subfigure[Case$2S2$]
        {\label{fig:speciesGladstoneDale2S2}\includegraphics[width=0.48\textwidth]{chapter4/fiveVsElevenSpecies/2S2_speciesGladstoneDale}}
        \caption{\label{fig:speciesGladstoneDaleStudy}Species mass
        fraction results for cases in \Cref{tab:speciesStudyCases}}
    \end{figure}
    %% Species Gladstone-Dale %%

    %% Total Gladstone-Dale %%
    \Cref{fig:totalGladstoneDaleStudy} shows the total \gls{GD} for the cases 
    described in \Cref{tab:speciesStudyCases} and \Cref{tab:speciesStudyMassFraction}. 
    As expected from \Cref{fig:speciesGladstoneDaleStudy}, the contribution
    from the ion species is minimal and not noticeable in the total \gls{GD}.
    \begin{figure}[H]
        \centering
        \subfigure[Case$1S1$]
        {\label{fig:totalGladstoneDale1S1}\includegraphics[width=0.45\textwidth]{chapter4/fiveVsElevenSpecies/1S1_totalGladstoneDale}}
        \subfigure[Case$1S2$]
        {\label{fig:totalGladstoneDale1S2}\includegraphics[width=0.45\textwidth]{chapter4/fiveVsElevenSpecies/1S2_totalGladstoneDale}}
        \subfigure[Case$2S1$]
        {\label{fig:totalGladstoneDale2S1}\includegraphics[width=0.45\textwidth]{chapter4/fiveVsElevenSpecies/2S1_totalGladstoneDale}}
        \subfigure[Case$2S2$]
        {\label{fig:totalGladstoneDale2S2}\includegraphics[width=0.45\textwidth]{chapter4/fiveVsElevenSpecies/2S2_totalGladstoneDale}}
        \caption{\label{fig:totalGladstoneDaleStudy}Total \gls{GD}
        results for cases in \Cref{tab:speciesStudyCases}}
    \end{figure}
    %% Total Gladstone-Dale %%

    %% Index of Refraction %% 
    \Cref{fig:refractionIndexStudy} shows the total index of refraction for the
    cases described in \Cref{tab:speciesStudyCases} and
    \Cref{tab:speciesStudyMassFraction}. As expected from
    \Cref{fig:refractionIndexStudy}, the contribution from the ion species is
    minimal and not noticeable in the index of refraction. Hence, it can be
    concluded that using a five-species gas is sufficient to predict the index
    of refraction, making the computational cost less expensive.
    \begin{figure}[H]
        \centering
        \subfigure[Case$1S1$]
        {\label{fig:refractionIndex1S1}\includegraphics[width=0.45\textwidth]{chapter4/fiveVsElevenSpecies/1S1_refractionIndex}}
        \subfigure[Case$1S2$]
        {\label{fig:refractionIndex1S2}\includegraphics[width=0.45\textwidth]{chapter4/fiveVsElevenSpecies/1S2_refractionIndex}}
        \subfigure[Case$2S1$]
        {\label{fig:refractionIndex2S1}\includegraphics[width=0.45\textwidth]{chapter4/fiveVsElevenSpecies/2S1_refractionIndex}}
        \subfigure[Case$2S2$]
        {\label{fig:refractionIndex2S2}\includegraphics[width=0.45\textwidth]{chapter4/fiveVsElevenSpecies/2S2_refractionIndex}}
        \caption{\label{fig:refractionIndexStudy}Index of Refraction
        results for cases in \Cref{tab:speciesStudyCases}}
    \end{figure}
    %% Index of Refraction %% 

    %% Electric Susceptibility %% 
    \Cref{fig:electricSusceptibilityStudySpecies} shows the electric
    susceptibility for the cases described in \Cref{tab:speciesStudyCases} and
    \Cref{tab:speciesStudyMassFraction}. As expected from
    \Cref{fig:electricSusceptibilityStudySpecies}, the contribution from the ion species is
    minimal and not noticeable in the electric susceptibility of refraction. Hence, it can be
    concluded that using a five-species gas is sufficient to predict the
    electric susceptibility, making the computational cost less expensive.
    \begin{figure}[H]
        \centering
        \subfigure[Case$1S1$]
        {\label{fig:electricSusceptibility1S1}\includegraphics[width=0.45\textwidth]{chapter4/fiveVsElevenSpecies/1S1_electricSusceptibility}}
        \subfigure[Case$1S2$]
        {\label{fig:electricSusceptibility1S2}\includegraphics[width=0.45\textwidth]{chapter4/fiveVsElevenSpecies/1S2_electricSusceptibility}}
        \subfigure[Case$2S1$]
        {\label{fig:electricSusceptibility2S1}\includegraphics[width=0.45\textwidth]{chapter4/fiveVsElevenSpecies/2S1_electricSusceptibility}}
        \subfigure[Case$2S2$]
        {\label{fig:electricSusceptibility2S2}\includegraphics[width=0.45\textwidth]{chapter4/fiveVsElevenSpecies/2S2_electricSusceptibility}}
        \caption{\label{fig:electricSusceptibilityStudySpecies}Electric susceptibility
        results for cases in \Cref{tab:speciesStudyCases}}
    \end{figure}
    %% Electric Susceptibility %% 
